% Transcription — fonds Alexandre Grothendieck, shelfmark 137, batch 2 (pages 21-40).
% Established from the facsimile archives/batches/137/batch-02.pdf (a copy of
% 137.pdf fetched through the project's relay,
% https://grothendieck-relay.onrender.com/source/137.pdf; PDF page k is
% archivists' page 20+k, no cover sheet), per the skill transcribe-grothendieck.
% The folder is eighty-five pages in five batches (1-20, 21-40, 41-60, 61-80,
% 81-85); this is the second. Batches 1 and 3 were read after drafting and the
% notation aligned with them (Lambda^i Phi_*, Lambda^j Psi_*, Phi(i,n) as in
% batch 1; C, A, C^i_*, K, C^•, Ab_k as in batch 3).
%
% Pass: Opus 5.5 (claude-opus-5-5), 2026-09-26 — first pass, unchecked against the pages by a human.
%
% Pages skipped: none. Pages summarised: none. Page 22 is a double-width sheet
% written in two columns and is given column by column. The upper half of
% page 40 is cancelled by three long diagonals and is transcribed in full; the
% lower half of page 38 is blank. Pencil throughout; pages 22-30 and 33-37 on
% music-ruled paper, the rest on hole-punched lined sheets. The formulas carry;
% much connective prose on pages 36-37 does not and is left \ill{} or
% \uncertain{}.
%
% His pagination: none on these leaves. His numbered runs: on p. 22 the
% conditions (a), (b), (c) (circled) and d), with a second circled (c) on
% p. 23 (the hypothesis that e_i is a basis of Ext^1); a circled « d » in the
% margin of p. 26; « Prop. » (unnumbered) on p. 31; 1°) and 2°) for the two
% cases C (abelian) / C_0 (additive) on p. 36, and 1°) / 2°) again for the
% two programmes (Dedekind ring / arbitrary k) on p. 37; on p. 39 the
% conditions a), b) on (K, C^•) and 1°)-3°) for complexes of chains; on p. 40
% a), b), « Prop 1 » and « Exemple 1 » — the « prop. 1 », « Ex 2 » of batch 3
% (p. 41) continue this run. All recorded in \note{}s where they stand.
%
% What the batch is about. Pp. 21-27 end the Hom computations begun on p. 20
% (batch 1) for strictly ∂-perfect semi-simplicial modules over a principal
% ring: Hom between Lambda^i Psi_*, Lambda^j Phi_* and the indecomposables
% Phi(i,n), the extension Phi(i,n) obtained from Lambda^i Phi_* by
% n·id on Lambda^{i-1} Psi_*; p. 22 restates everything axiomatically in the
% additive category C_0 with a cochain complex Phi^• and Psi^i = Z(Phi^•),
% and pp. 23-26 compute Hom(Phi(i,n), Phi(j,m)) (= k for j = i, via
% qm = pn and m' = m/(m,n)), Hom(Phi(i,-), Psi^j), Hom(Psi^i, Phi(j,-));
% p. 27 concludes that the category is known from the full subcategory A of
% these objects with its « ½ additive » structure, and names the goal: the
% ⊗-structure and the operations Gamma^i. Pp. 28-35: general reconstruction of
% a category C from a full subcategory A — C -> Â is fully faithful when every
% object is a colimit of objects of A, the essential image when C has limits
% and colimits of given types T, T' (Prop. p. 31), the equivalence
% Â_T ≅ (°_T')° (duality of free modules as the one-object case),
% Hom_add(C, C') ≅ Hom_½add(A, C'), and the pairing F ⊗_C G / Hom_C(F, G).
% Pp. 36-38: the programme — S = Δ^ = Ss, the functor
% M_• -> Hom(X_•, M_•) on ∂-perfect groups, and the aim of recovering the
% cochains C^*(X_•, k) with their multiplicative structure (⊗ and Gamma^i,
% divided powers), De Rham-Sullivan complexes, the homotopy type of X_•
% (« faisceautiser et relativiser... !!! »). Pp. 39-40: the axiomatic setting
% of batch 3 — a k-additive K with a complex C^• satisfying a) Hom(C^i, C^j)
% and b), realised by acyclic bounded complexes of free modules, and the
% category C of flasque str. ∂-perfect semi-simplicial modules, sums of
% C^i_* = Lambda^{i+1} Phi_{*k}, with Prop 1 (k-linear functors out of C are
% cochain complexes of B) and Exemple 1.
%
% Notation fixed once for the batch: \Lambda^{i}\Psi_{*}, \Lambda^{i}\Phi_{*},
% \Phi(i,n) (also \Phi_{\bullet}(i,n) where he writes it so) as in batch 1;
% \Phi^{i}, \Psi^{i}, \partial_i, u_i, e_i on pp. 22-26; C, A, \widehat{A},
% \widehat{A}_{ab}, T, T' on pp. 28-35 (script T for his curly T); \mathcal{S}
% for his S; \mathbb{K}, \mathbb{C}^{\bullet} on p. 39, \mathcal{C},
% \mathcal{A}, C^{i}_{*} on p. 40, following batch 3. His underlined lim is
% set \varinjlim or \varprojlim by the arrow he draws; « ½ add » as
% \frac{1}{2}add in indices.
%
% Readings to check first: p. 21 last sentence (« cat. »); p. 22 the
% exponents of the cochain complex and the small marginal diagram; p. 23 the
% letters m, n, m', n' (he appears to swap them); p. 24 the bottom row of the
% diagram (indices overwritten) and the whole paragraph under it; p. 26 the
% class « d e_n = d u e_i »; p. 27 first paragraph; pp. 36-37 nearly all
% connective prose; p. 40 the rotated margin note (« On suppose k ... »).
\documentclass[11pt,a4paper]{article}
% Le préambule commun à toutes les transcriptions du fonds Grothendieck.
%
% Il tient en une idée : **l'incertitude se compose, elle ne se résout pas.**
% Une transcription qui lisse un mot douteux a détruit la seule chose pour
% laquelle on la faisait — un lecteur doit pouvoir savoir, à chaque endroit,
% ce qui était sur le feuillet et ce qui a été ajouté. Les six macros qui
% suivent sont donc l'appareil critique, et non de la décoration.
%
% Les mêmes commandes sont reconnues par `scripts/render.mjs`, qui produit la
% vue de lecture du site. Une macro ajoutée ici doit l'être aussi là-bas,
% faute de quoi le rendu échouera bruyamment — ce qui est voulu : un rendu
% silencieusement faux est pire qu'un rendu absent.

\ProvidesPackage{grothendieck}[2026/08/08 Transcriptions du fonds Grothendieck]

\usepackage{amsmath,amssymb,amsthm}
\usepackage{mathrsfs}
\usepackage{stmaryrd}
% Les diagrammes commutatifs sont partout dans ce fonds : c'est la langue
% même de la géométrie algébrique de Grothendieck, et une transcription qui
% les rendrait en prose ne transcrirait rien.
\usepackage{tikz-cd}
% Français par défaut — c'est la langue des feuillets — mais l'anglais est
% chargé aussi : la traduction et le résumé sont des documents anglais, et la
% typographie française (espaces avant les deux-points) y serait une faute.
% Ces documents basculent par \selectlanguage{english} en tête de corps.
\usepackage[english,french]{babel}
\usepackage{fontspec}
\usepackage{geometry}
\geometry{a4paper,margin=25mm,marginparwidth=28mm}
\usepackage{marginnote}
\usepackage{xcolor}
% ulem, not soul. `soul`'s \st and \ul are built on a character-by-character
% scan that XeLaTeX's font machinery breaks: the document fails with an
% undefined control sequence rather than degrading. ulem does the same two jobs
% and is Unicode-engine safe. `normalem` stops it hijacking \emph, which the
% transcriptions use for Grothendieck's own emphasis.
\usepackage[normalem]{ulem}
\usepackage{hyperref}
\hypersetup{colorlinks=true,linkcolor=black,urlcolor=black,pdfborder={0 0 0}}

% --- Métadonnées du lot -----------------------------------------------------
% Reprises telles quelles par le script de rendu, qui les affiche en tête de
% la vue de lecture. Elles fixent ce que le fichier prétend couvrir : sans
% elles, une transcription flotte sans point d'attache dans un fonds de seize
% mille feuillets.

\newcommand{\folder}[1]{\gdef\grothFolder{#1}}
\newcommand{\batch}[1]{\gdef\grothBatch{#1}}
\newcommand{\leaves}[2]{\gdef\grothFirst{#1}\gdef\grothLast{#2}}
\newcommand{\foldertitle}[1]{\gdef\grothFoldertitle{#1}}
\gdef\grothFolder{?}\gdef\grothBatch{?}\gdef\grothFirst{?}\gdef\grothLast{?}\gdef\grothFoldertitle{}

% --- Le résumé --------------------------------------------------------------
%
% Il ouvre la lecture modernisée, sous le titre « Résumé », et il est composé
% en petit. Ce n'est pas un ornement : c'est le seul passage du document qui
% ne soit pas des mathématiques, et un lecteur doit pouvoir le reconnaître
% avant de l'avoir lu — puis le sauter s'il connaît déjà le sujet, ou s'y
% arrêter s'il ne le connaît pas. Le filet qui le suit marque la reprise du
% corps.

\newenvironment{resume}
  {\small\setlength{\parskip}{0.45em}}
  {\par\medskip\hrule\medskip}

% --- Mots-clés ---------------------------------------------------------------
%
% En anglais, en clôture du résumé de la lecture modernisée : le vocabulaire
% moderne sous lequel chercher ce que les feuillets construisent. Le manifeste
% du site (`npm run manifest`) les extrait pour étiqueter la cote entière —
% cette ligne est la source unique des tags, il n'y a pas de fichier à côté.
\newcommand{\keywords}[1]{\par\smallskip\noindent
  {\footnotesize\textsc{Keywords} --- \textit{#1}}\par}

% --- L'appareil critique ----------------------------------------------------

% Le numéro de feuillet, en marge. C'est l'ancre qui permet de régler un
% désaccord sur une formule en regardant *un* feuillet plutôt qu'une chemise
% de six cents. Le site s'en sert aussi pour tourner les pages du fac-similé
% quand on fait défiler la transcription.
\newcommand{\leaf}[1]{%
  \marginnote{\footnotesize\color{gray!70}\textsf{#1}}%
  \label{leaf:#1}\ignorespaces}

% Illisible : on ne devine pas. Un mot inventé qui se lit comme les autres est
% le pire résultat possible.
\newcommand{\ill}{\textcolor{red!60!black}{[\,\ldots\,]}}

% Lecture douteuse : le mot est donné, et signalé comme incertain.
\newcommand{\uncertain}[1]{\uline{#1}}

% Ajout de l'éditeur — une lettre manquante, un mot sous-entendu.
\newcommand{\add}[1]{\textcolor{blue!50!black}{[#1]}}

% Biffé par Grothendieck lui-même. Ce qu'il a rayé fait partie du document :
% une rature dit souvent où la pensée a changé de direction.
\newcommand{\struck}[1]{\sout{#1}}

% Note du transcripteur, sur l'état matériel du feuillet.
\newcommand{\note}[1]{\par\smallskip{\small\color{gray!60!black}\textit{[#1]}}\par\smallskip}

% Note marginale de Grothendieck, distincte du corps du texte.
\newcommand{\marginal}[1]{\par\smallskip\noindent\hspace{1em}%
  {\small\color{gray!60!black}#1}\par\smallskip}

% --- Titre ------------------------------------------------------------------

\AtBeginDocument{%
  \begin{center}
    {\large\bfseries Fonds Alexandre Grothendieck --- cote n\textsuperscript{o}~\grothFolder}\\[2pt]
    {\itshape\grothFoldertitle}\\[4pt]
    {\small feuillets \grothFirst--\grothLast\ (lot \grothBatch)}\\[2pt]
    {\footnotesize\color{gray!60!black}Transcription établie d'après le fac-similé numérisé
     par l'Université de Montpellier.\\
     Les lectures incertaines sont soulignées, les illisibles notées [\,\ldots\,],
     les ajouts entre crochets.}
  \end{center}
  \vspace{1em}}

\endinput


\folder{137}
\batch{2}
\pages{21}{40}
\dating{[vers 1975-1976]}
\watermark{Édition de démonstration}
\foldertitle{Dévissage des complexes ½ simpliciaux ∂ - parfaits et structures multiplicatives… (1975 ou 1976) : notes manuscrites (s.d.)}

\begin{document}

\page{21}
\[
\operatorname{Hom}_{k}(\Lambda^{i}\Psi_{*}, \Lambda^{j}\Psi_{*}) \simeq
\begin{cases}
0 & \text{si } j \neq i \\
k & \text{si } j = i
\end{cases}
\]
de façon précise
\[
\operatorname{Hom}_{k}(\Lambda^{i}\Psi_{*}, \Lambda^{i}\Psi_{*}) = k\,\mathrm{id}_{\Lambda^{i}\Psi_{*}} \simeq k ,
\]
et
\[
\operatorname{Hom}_{k}(\Lambda^{i}\Psi_{*}, \Phi_{\bullet}(j,n)) \simeq
\begin{cases}
0 & \text{si } j \neq i+1 \\
k & \text{si } j = i+1
\end{cases}
\]
\note{dans la première ligne de l'accolade, un symbole biffé avant « $i+1$ »}
en particulier
\[
\operatorname{Hom}_{k}(\Lambda^{i}\Psi_{*}, \Lambda^{j}\Phi_{*}) \simeq
\begin{cases}
0 & \text{si } j \neq i+1 \\
k & \text{si } j = i+1
\end{cases}
\]
\note{l'exposant du second argument est surchargé ($j$ corrigé ?)}
et
\[
\operatorname{Hom}_{k}(\Lambda^{i}\Psi_{*}, \Lambda^{i+1}\Phi_{*}) = k\,\partial \simeq k
\]
où $\partial : \Lambda^{i}\Psi_{*} \hookrightarrow \Lambda^{i+1}\Phi_{*}$ est l'inclusion canonique.

\emph{NB} On a une \add{homom.\ de} suites exactes \uncertain{canoniques}

\begin{tikzcd}[column sep=small, row sep=small]
0 \arrow[r] & \Lambda^{i-1}\Psi_{*} \arrow[r] \arrow[d, "n\,\mathrm{id}"] & \Lambda^{i}\Phi_{*} \arrow[r] \arrow[d] & \Lambda^{i}\Psi_{*} \arrow[r] \arrow[d, no head, "\mathrm{id}"] & 0 \\
0 \arrow[r] & \Lambda^{i-1}\Psi_{*} \arrow[r] & \Phi_{\bullet}(i,n) \arrow[r] & \Lambda^{i}\Psi_{*} \arrow[r] & 0
\end{tikzcd}

\note{la flèche verticale médiane porte « $=$ » biffé et une flèche ; celle de droite est doublée, marquée
« id »}

qui montre que \uncertain{$\Phi_{\bullet}(i,n) \to$ \ill{} \uncertain{cat.}} \uncertain{se} déduit de
\add{\uncertain{\ill{} cat.}} $\Lambda^{i}\Phi_{*}$ de $\Lambda^{i}\Psi_{*}$ par $\Lambda^{i-1}\Psi_{*}$ par
l'hom.\ $n\,\mathrm{id} : \Lambda^{i-1}\Psi_{*} \to \Lambda^{i-1}\Psi_{*}$.\note{passage rapide ; « cat. » est
peut-être l'abréviation de « extension canonique »}

\page{22}
\note{feuillet double, écrit en deux colonnes ; la colonne de gauche est donnée d'abord}

\emph{En résumé.} Dans la catégorie \add{$k$-lin.\ $\mathcal{C}$} \struck{abélienne} de tous les $k$-modules ½
simpliciaux, ceux qui sont \add{str.\ $\partial$-}parfaits forment une sous-catégorie \emph{additive}
$\mathcal{C}_0$ \struck{\ill{}}, dans laquelle on a un complexe de cochaînes
\[
0 \to \Phi^{0} \to \Phi^{1} \to \Phi^{2} \to \cdots \to \Phi^{i} \xrightarrow{\ \partial_i\ } \Phi^{i+1} \to \cdots
\]
\note{chaque exposant est récrit sur une première lecture, peut-être d'autres indices ; les
$\Phi^{i} = \Lambda^{i+1}\Phi_{*}$ du lot 1}
satisfaisant à ceci :
\begin{enumerate}
\item[(a)] \uncertain{Complexe} tot.\ acyclique dans $\mathcal{C}$, \struck{\ill{}}
\item[(b)] Pour tout $i, j \geq 0$, on a
\[
\operatorname{Hom}(\Phi^{i}, \Phi^{j}) \simeq
\begin{cases}
0 & \text{si } j \neq i, i+1 \\
\simeq k \ (\text{engendré par } \mathrm{id}) & \text{si } j = i \\
\simeq k \ (\text{engendré par } \partial_i) & \text{si } j = i+1
\end{cases}
\]
\item[(c)] Posant $\Psi^{i} = Z^{i+1}(\Phi^{\bullet}) = \operatorname{Ker}(\Phi^{i+1} \xrightarrow{\partial_{i+1}}
\Phi^{i+2}) \simeq \operatorname{Coker}_{\mathcal{C}}(\Phi^{i-1} \xrightarrow{\partial_{i-1}} \Phi^{i})$,
\struck{on a donc $\operatorname{Hom}(\Psi^{i}$}
\end{enumerate}
\note{(a), (b), (c) sont cerclés ; le (c) est surchargé. Les indices de $Z$ et de $\operatorname{Ker}$ sont
d'une lecture douteuse. Dans la marge gauche, en regard de (c) : « \uncertain{unicité} \uncertain{de}
\uncertain{(c)} »}
\[
\Psi^{0} \simeq \Phi^{0}
\]
\[
(*) \qquad 0 \to \Psi^{i-1} \xrightarrow{\ \partial_i\ } \Phi^{i} \xrightarrow{\ u_i\ } \Psi^{i} \to 0
\qquad \text{ext.\ can.}
\]
\note{sous cette suite, un croquis : $\Phi^{i-1} \xrightarrow{u} \Psi^{i-1}$ (flèche montante), des obliques
$\partial_{i-1}$ vers $\Phi^{i}$ et $\partial_i$ de $\Phi^{i}$ vers $\Phi^{i+1}$, et une flèche descendante
$\Psi^{i} \hookrightarrow \Phi^{i+1}$}
\[
\begin{aligned}
&\operatorname{Hom}(\Psi^{i}, \Psi^{j}) \ (\subset \operatorname{Hom}(\Phi^{i}, \Phi^{j+1})) \simeq
\begin{cases}
0 & \text{si } j \neq i \\
k & \text{si } j = i \ (\text{base } \mathrm{id})
\end{cases} \\
&\operatorname{Hom}(\Phi^{i}, \Psi^{j}) \ (\subset \operatorname{Hom}(\Phi^{i}, \Phi^{j+1})) \simeq
\begin{cases}
0 & \text{si } j \neq i \\
k & \text{si } j = i \ (\text{base } u_i)
\end{cases} \\
&\operatorname{Hom}(\Psi^{i}, \Phi^{j}) \ (\subset \operatorname{Hom}(\Phi^{i}, \Phi^{j})) \simeq
\begin{cases}
0 & \text{si } j \neq i+1 \ (i \neq 0) \\
k & \text{si } j = i+1 \ (\text{base } \partial_i) \\
k & \text{si } j = i = 0 \ (\text{base } \mathrm{id})
\end{cases}
\end{aligned}
\]
\note{dans la première accolade, le second membre de « $j \neq i$ » est surchargé ; dans la deuxième, on lit
« $j = \theta$ » pour $j = i$}

d) Soit, \struck{\ill{}} pour $n \in k - 0$, \add{et $i \geq 1$}, $\Phi(i,n)$ l'ext.\ de $\Psi^{i}$ par
$\Psi^{i-1}$ déduite de l'ext.\ $\Phi^{i}$ $(*)$ par $n\,\mathrm{id}_{\Psi^{i-1}}$, de sorte qu'on a un homom.\
de suites ex.\ \uncertain{courtes}

\begin{tikzcd}[column sep=small, row sep=small]
0 \arrow[r] & \Psi^{i-1} \arrow[r, "\partial_i"] \arrow[d, "n\,\mathrm{id}"] & \Phi^{i} \arrow[r, "u_i"] \arrow[d] & \Psi^{i} \arrow[r] \arrow[d, "\mathrm{id}"] & 0 \\
0 \arrow[r] & \Psi^{i-1} \arrow[r] & \Phi(i,n) \arrow[r] & \Psi^{i} \arrow[r] & 0
\end{tikzcd}

(\emph{NB} $\Phi(i,1) = \Phi^{i}$.)\note{ce diagramme ouvre la colonne de droite}
\marginal{on a aussi
$0 \to \Psi^{i-1} \to \Phi(i,n) \to \Psi^{i} \to 0$ au-dessus de
$0 \to \Psi^{i-1} \to \Phi^{i} \to \Psi^{i} \to 0$, flèches $n\,\mathrm{id}$, $\rho$, $\mathrm{id}$}\note{écrit en
petit en haut à droite, relié par un trait oblique au $\operatorname{Hom}$ qui suit ; lecture douteuse}
Alors on a
\[
\operatorname{Hom}(\Phi(i,-), \Phi(j,-)) \to \operatorname{Hom}(\Psi^{i-1}, \Psi^{j})
\]
et si $j \neq i+1$, ce dernier Hom est nul, \uncertain{donc} si $j \neq i+1$, la structure d'extension est
respectée par \uncertain{tout} homom.\ $\Phi(i,-) \to \Phi(j,-)$

\begin{tikzcd}[column sep=small, row sep=small]
0 \arrow[r] & \Psi^{i-1} \arrow[r] \arrow[d] & \Phi(i,-) \arrow[r] \arrow[d] & \Psi^{i} \arrow[r] \arrow[d] & 0 \\
0 \arrow[r] & \Psi^{j-1} \arrow[r] & \Phi(j,-) \arrow[r] & \Psi^{j} \arrow[r] & 0
\end{tikzcd}

de plus si $j \neq i$, les flèches \uncertain{extrêmes} \uncertain{étant} nulles, donc on trouve
\[
\operatorname{Hom}(\Phi(i,n), \Phi(j,m)) \xleftarrow{\ \approx\ } \operatorname{Hom}(\Psi^{i}, \Psi^{j-1})
\simeq
\begin{cases}
k & \text{cas } j = i+1 \\
0 & \text{si } j \neq i+1
\end{cases}
\]
\marginal{\uncertain{engendré} \uncertain{par} \uncertain{le} \uncertain{comp.}\
$\Phi(i,n) \to \Psi^{i} \to \Phi(i+1,m)$}\note{note oblique entourée dans la marge gauche de la colonne, reliée
au cas $j = i+1$}

\emph{Cas $j = i$.} Tout homom.\ \uncertain{induit} \uncertain{un} $p\,\mathrm{id}$ sur $\Psi^{i-1}$ et
$q\,\mathrm{id}$ sur $\Psi^{i}$,

\begin{tikzcd}[column sep=small, row sep=small]
0 \arrow[r] & \Psi^{i-1} \arrow[r] \arrow[d, "p\,\mathrm{id}"] & \Phi(i,n) \arrow[r] \arrow[d, "\alpha"] & \Psi^{i} \arrow[r] \arrow[d, "q\,\mathrm{id}"] & 0 \\
0 \arrow[r] & \Psi^{i-1} \arrow[r] & \Phi(i,m) \arrow[r] & \Psi^{i} \arrow[r] & 0
\end{tikzcd}

\uncertain{déterminés} \uncertain{par} \uncertain{définition} par $\alpha$. On a alors (puisque \uncertain{par
définition} $\alpha$ \uncertain{est} défini \uncertain{par} $p$ \uncertain{et} $q$) que l'élément de
$\operatorname{Ext}^{1}(\Psi^{i}, \Psi^{i-1})$ défini par $q^{*}(e_{i,m}) - p_{*}(e_{i,n})$ (où $e_{i,m}$,
$e_{i,n}$ \uncertain{définissent} \uncertain{resp.} les ext.\ $\Phi(i,-)$, $\Phi(i,m)$)
\note{le bas de la colonne est très serré ; la suite est en page 23}

\page{23}
est nul, or
\[
q^{*}(e_{i,m}) = m\,e_i , \qquad p_{*}(e_{i,n}) = n\,e_i
\]
\note{les deux indices se lisent « $m$ » sous surcharge ; le second est vraisemblablement $n$, que demande la
suite}
où $e_i$ définit l'extension $\Phi^{i}$. Donc la condition est que
\[
(q\,m - p\,n)\,e_i = 0 .
\]
On fait l'hypothèse
\[
\text{(c)} \qquad e_i \text{ est une base de } \operatorname{Ext}^{1}_{\mathcal{C}_0}(\Psi^{i}, \Psi^{i-1})
\]
\note{hypothèse encadrée, avec « (c) » cerclé en regard. Il y a donc deux (c) : celui de la page 22 et
celui-ci}
de sorte que la condition devient
\[
q\,m = p\,n
\]
i.e.
\[
\frac{q}{p} = \frac{n}{m} = \frac{n'}{m'}
\]
si $m' = \dfrac{m}{(m,n)}$, $n' = \dfrac{n}{(m,n)}$, donc on aura
\[
p = \struck{\ill{}}\ t\,m' , \qquad q = t\,n' \qquad (t \in k)
\]
\note{on lit « $q = t\,m'$ » ; la relation $q/p = n'/m'$ demande $t\,n'$, que nous gardons. Même glissement
possible des lettres, sur la page, entre $m$ et $n$}
Or pour $p, q$ correspondant à un $t$ (\uncertain{d'ailleurs} unique) \uncertain{dès} \uncertain{que}
\uncertain{déterminés} par $p, q$ l'hom.\ $\alpha$ qui leur correspond est déterminé mod.\ un
$\operatorname{Hom}(\Psi^{i}, \Psi^{i-1}) = 0$, donc est unique. On trouve donc
\[
\operatorname{Hom}(\Phi(i,n), \Phi(i,m)) \simeq k .
\]

\page{24}
Car \uncertain{$j = i-1$}\note{le $i-1$ est récrit sur un autre indice (peut-être $i+1$)} :

\begin{tikzcd}[column sep=small, row sep=small]
 & \Psi^{i-1} \arrow[r] \arrow[d, "u\circ d"] & \Phi^{i} \arrow[r] \arrow[d] & \Psi^{i} \arrow[d, no head, "="] & \\
0 \arrow[r] & \Psi^{i-1} \arrow[r, "D"] \arrow[d, "\wr"] & \Phi(i,-) \arrow[r] \arrow[d, "\alpha=?"] & \Psi^{i} \arrow[r] \arrow[d, "\wr"] & 0 \\
0 \arrow[r] & \Psi^{i-1} \arrow[r] & \Phi(i,m) \arrow[r, "u"] & \Psi^{i} \arrow[r] & 0
\end{tikzcd}

\note{la ligne du bas est surchargée : les exposants de $\Psi$ ont été récrits (on lit sous la correction $i+2$ à gauche et un indice biffé à droite), et le second argument de $\Phi$ se lit \uncertain{$m$}, le premier peut-être $i-1$ corrigé en $i$ ; nous donnons la lecture qui s'accorde avec la suite de la page, où $m$ et $n$ interviennent}

Le composé $\Phi^{i} \to \Phi(i,-) \to \Phi(i-1,-) \to \Psi^{i}$\note{lecture de la chaîne incertaine} est nul car
$\underline{\mathrm{Hom}}(\Phi^{i}, \Psi^{i-1}) = 0$ ; \uncertain{$\theta = u\alpha D$} : donc l'hom.\
\struck{$\Psi^{i} \to \Psi^{i}$} $\Psi^{i-1} \to \Psi^{i-1}$ composé \ill{} \ill{} \uncertain{id est nul}, i.e.\
\uncertain{$u\theta = 0$}, or $\underline{\mathrm{Hom}}(\Psi^{i-1}, \Psi^{i-1}) \simeq k$, donc $\theta = 0$, donc
$\alpha$ respecte la structure d'extension, et on \uncertain{fait comme} plus haut pour conclure
$\alpha = 0$. Donc on retrouve
\[
\underline{\mathrm{Hom}}(\Phi(i,-), \Phi(j,-)) \simeq
\begin{cases}
0 & \text{si } j \neq i, i+1 \\
k & \text{si } j = i+1 \\
k & \text{si } j = i
\end{cases}
\]
\note{en regard des deux dernières lignes de l'accolade : pour $j = i+1$, « (base \uncertain{en consid.} l'hom.\ can.\
$\Phi(i,-) \to \Psi^{i} \to \Phi(i+1,-)$) » ; pour $j = i$, « (\uncertain{engendré par} \ill{} respectant la
struct.\ d'ext.) »}

\uncertain{et étant} déterminé par ce qu'il induit sur $\Psi^{i}$ et $\Psi^{i-1}$, qui sont \uncertain{deux
éléments} de la forme \uncertain{$t\,m'\,\mathrm{id}_{\Psi^{i-1}}$ et $t\,n'\,\mathrm{id}_{\Psi^{i}}$}), où
\[
n' = \frac{n}{(m,n)}, \qquad m' = \frac{m}{(m,n)}
\]
(il y a une base \uncertain{can.} une fois qu'on a choisi le \uncertain{pgcd} $(m,n)$ de $m$ et $n$).

\page{25}
\[
\underline{\mathrm{Hom}}(\Phi(i,-), \Psi^{j}) \subset \underline{\mathrm{Hom}}(\Phi(i,-), \Phi^{j+1})
\]
\note{sous $\Phi^{j+1}$, un signe $\wr$ et $\Phi(j+1,-)$}
$= 0$ si $j+1 \neq i, i+1$, i.e.\ si $j \neq i-1, i$ ;
\[
\underline{\mathrm{Hom}}(\Psi^{i-1}, \Psi^{j}) = 0 \quad \text{si } j \neq i-1 .
\]
\note{une flèche verticale descend du premier $\underline{\mathrm{Hom}}$ vers cette ligne}

Si donc $j = i$, $\underline{\mathrm{Hom}}(\Phi(i,-), \Psi^{i}) \hookrightarrow \underline{\mathrm{Hom}}(\Phi^{i}, \Psi^{i})
\simeq k$ (base \uncertain{l'hom.\ can.}).

Si $j = i-1$ :

\begin{tikzcd}[column sep=small, row sep=small]
 & \Psi^{i-1} \arrow[r] \arrow[d, "u\circ d"] & \Phi^{i}(k) \arrow[r] \arrow[d] & \Psi^{i} \arrow[r] \arrow[d] & 0 \\
0 \arrow[r] & \Psi^{i-1} \arrow[r, "i"] & \Phi(i,-) \arrow[r] \arrow[d, "\alpha=?"] & \Psi^{i} \arrow[r] & 0 \\
 & & \Psi^{i-1} & &
\end{tikzcd}

\note{l'argument de $\Phi^{i}$ dans la ligne du haut est surchargé ; « $(k)$ » est une lecture douteuse}

Le composé de $\alpha i$ avec \uncertain{$\Phi^{i} \to \Phi(i,-)$}\note{sous la flèche, un mot biffé
(\uncertain{$u\circ d$}) et $\Psi^{i}$} est dans $\underline{\mathrm{Hom}}(\Phi^{i}, \Psi^{i-1}) = 0$, donc
\struck{$u\alpha i \in \underline{\mathrm{Hom}}$} $\alpha i \circ u\,\mathrm{id}_{\Psi^{i-1}} = u(\alpha i)$
\uncertain{$\in \underline{\mathrm{Hom}}(\Psi^{i-1}, \Psi^{i-1})$} est nul, or comme
$\underline{\mathrm{Hom}}(\Psi^{i-1}, \Psi^{i-1}) \simeq k$, on voit que $\alpha i = 0$, donc $\alpha$ provient
\uncertain{d'une flèche} \ill{} par $\Psi^{i} \to \Psi^{i-1}$, donc est nul. Donc
\[
\underline{\mathrm{Hom}}(\Phi(i,-), \Psi^{j}) \simeq
\begin{cases}
0 & \text{si } j \neq i \\
\simeq k & \text{si } j = i
\end{cases}
\]
\note{en regard de la seconde ligne : « (base l'hom.\ can.\ $\Phi(i,-) \to \Psi^{i}$) »}

\page{26}
\[
\underline{\mathrm{Hom}}(\Psi^{i}, \Phi(j,-)) \longrightarrow \underline{\mathrm{Hom}}(\Psi^{i}, \Psi^{j})
\]
\note{sous le second membre, « $\wr$ » puis « 0 si $j \neq i$ » ; sous le premier, un mot \uncertain{« \ill{} »} en petits caractères}
donc si $j \neq i$, $\underline{\mathrm{Hom}}(\Psi^{i}, \Phi(j,-)) \simeq
\underline{\mathrm{Hom}}(\Psi^{i}, \Psi^{j-1})$ \struck{\ill{}} qui est nul si $i \neq j-1$ i.e.\
$j \neq i+1$ ; si $j = i+1$, on trouve $\simeq k$.

Pour $j = i$, le composé $\Psi^{i} \to \Phi(i,-) \xrightarrow{\text{can}} \Psi^{i}$ est de la forme
$d.\mathrm{id}_{\Psi^{i}}$, donc l'image \uncertain{inverse de l'extension}

\begin{tikzcd}[column sep=small, row sep=small]
 & & & \Psi^{i} \arrow[d, "\alpha"'] \arrow[dr] & \\
0 \arrow[r] & \Psi^{i-1} \arrow[r] & \Phi(j,-) \arrow[r] & \Psi^{i} \arrow[r] & 0
\end{tikzcd}

\note{le diagramme est à gauche de la page, dessiné plus librement : $\Psi^{i}$ en haut, une flèche $\alpha$ vers le
bas et une oblique ; nous le redessinons}

$\Phi(i,-)$ est de classe $d\,e_{n} = d\,u\,e_{i}$\note{indices douteux}, qui \uncertain{est nulle}, donc
$d\,e_{i} = 0$, donc $d = 0$, donc \struck{\ill{}} \uncertain{$\alpha$ se factorise} par
$\Psi^{i} \to \Psi^{i-1}$, donc est nul.

\struck{Pour $j = i+1$, \ldots\ $= \underline{\mathrm{Hom}}($}

Donc
\[
\underline{\mathrm{Hom}}(\Psi^{i}, \Phi(j,-)) \simeq
\begin{cases}
0 & \text{si } j \neq i+1 \\
k & \text{si } j = i+1
\end{cases}
\]
\note{en regard de la seconde ligne : « (base formée de l'hom.\ can.\ $\Psi^{i} \to \Phi(i+1,-)$) » ; avant la
formule, un $\underline{\mathrm{Hom}}($ biffé. Un trait ondulé sépare ensuite la page}

Ensuite \ill{} \uncertain{supposons} \ldots

\marginal{(d)}\note{« d » cerclé en tête de ligne, avec un trait vertical dans la marge gauche sur les lignes qui
suivent}
\ill{} \uncertain{de $C$} \ill{} \ill{} \ill{} de la forme finie \uncertain{d'objets} de la forme
$\Phi(i,n)$ $(n \geq 1)$ ou $\Psi^{j}$ $(j \geq 0)$.

\page{27}
En outre on sait que si $A$ est la sous-catégorie pleine de $C$ engendrée par les $\Phi(i,-)$ et les
$\Psi^{i}$, alors $C$ est connue à équivalence près (\uncertain{déf.} \ill{} \ill{} \ill{}) si on connaît $A$
avec sa structure ½ additive. Or c'est \add{elle} que \uncertain{nous venons} d'étudier \uncertain{(on a omis} de
décrire la composition des morphismes, qui est évidente\ldots).\note{« elle » est ajouté en interligne}
On trouve donc que $C \simeq$ \uncertain{groupes} ½ simpliciaux \uncertain{abéliens} \ill{} \ill{}
$\partial$-parfaits\ldots.

La chose \uncertain{intéressante} : faire \ill{} d'un \ill{} la $\otimes$-structure et les
\uncertain{opérations} \uncertain{$\Gamma^{i}$}.

\page{28}
Soit $C$ une catégorie, $A$ une sous-catégorie \add{pleine}.\note{« pleine » ajouté en interligne}
Considérons le foncteur composé
\[
C \xrightarrow{\text{can}} \widehat{C} \xrightarrow{\text{rest.}} \widehat{A} =
\underline{\mathrm{Hom}}(A^{\circ}, \mathrm{Ens})
\]
\note{un arc sous la flèche composée porte $\rho$}
si \struck{\ill{}} tout objet $X$ de $C$ est $\underline{\lim}$ \add{dans $C$} \struck{\ill{}} d'un diagramme de $A$
(\uncertain{ce qui} revient au même, $X \simeq \varinjlim_{Y \in \mathrm{Ob}(A/X)} Y$\note{la limite est
soulignée d'un trait simple, sans flèche lisible ; nous la rendons par $\varinjlim$, la limite canonique sur
$A/X$, et la variable est écrite au-dessous}), alors le foncteur $\rho$ est pl.\ fidèle \uncertain{(mais pas
nécess.\ $A$)}. Si $C$ est additive, le diagramme \add{ci-dessus} se complète en

\begin{tikzcd}[column sep=small, row sep=small]
C \arrow[r, "c_{ab}"] \arrow[dr, "c"'] \arrow[rr, bend left=30, "\rho_{ab}"] & \widehat{C}_{ab} \arrow[r, "\text{rest.}\,\rho_{ab}"] \arrow[d, "\text{oubli}\ \omega_C"] & \widehat{A}_{ab} \arrow[d, "\text{oubli}\ \omega_A"] \\
 & \widehat{C} \arrow[r, "\text{rest}\,\rho"] & \widehat{A}
\end{tikzcd}

\note{un grand arc passant sous le diagramme, de $C$ à $\widehat{A}$, porte $\rho$ ; nous l'omettons. Dans la
marge gauche, en oblique : « \uncertain{i.e.\ \ill{} structure} \ill{} \ill{} sur $C$ »}

avec $\omega_A$ fidèle, d'où il résulte que $\rho_{ab}$ est aussi pl.\ fidèle.
[C'est le cas en particulier si tout $X \in C$ est $\simeq$ somme d'objets de $C$.]\note{sic, on attendrait
« de $A$ »}

Supposons de plus que tout $X \in C$ est aussi \uncertain{somme} d'une $\varprojlim$ \add{dans $C$} d'objets de
$C$.\note{la $\lim$ porte ici une flèche vers la gauche} Alors l'image essentielle de $C$ dans $\widehat{A}$
(resp.\ $\widehat{A}_{ab}$) est contenue dans la sous-cat.\ pleine de $\widehat{A}$ resp.\ $\widehat{A}_{ab}$
formée des foncteurs

\page{29}
qui sont $\varprojlim$ du type envisagé de foncteurs représentables (dans $\widehat{A}$ ou $\widehat{A}_{ab}$
resp., \uncertain{pour} \ill{} dans $\widehat{A}_{ab}$ — la structure ½ additive de $A$ héritée par son
plongement dans $C$ additive). Lorsque de plus $C$ est stable par ce type de $\varprojlim$ (p.ex.\ des produits
finis), alors on a une \emph{équivalence} entre $C$ et la cat.\ des foncteurs $A^{\circ} \to \mathrm{Ens}$ (resp.\
$A^{\circ} \to \mathrm{Ab}$) qui sont de telles $\varprojlim$ de foncteurs représentables.
\struck{[Si p.ex.\ $A$ est réduit à un objet \ill{} dans $C$ au $C$ additive, \uncertain{resp.} si $A$ objet}
\note{ces lignes sont barrées d'un trait ondulé}
Si p.ex.\ $C$ est additive, donc $A$ induit sur $A$ une structure ½ additive, et si tout objet de $C$ est
isomorphe à une somme \add{finie} directe d'objets de $A$, alors $C \simeq$ sous-cat.\ pleine de
$\widehat{A}_{ab}$ (foncteurs additifs contravariants de $A$ dans $\mathrm{Ab}$) formée des foncteurs isom.\ à une
somme directe finie de foncteurs repr.

\page{30}
(Si $A$ est réduit à un objet \add{$A$ seul} d'anneau \uncertain{d'endomorphismes} $k$, on trouve une équivalence
avec la cat.\ des $k$-modules à droite libres de \emph{type fini}, par $X \mapsto \underline{\mathrm{Hom}}(A, X)$).

Transformons par dualité \uncertain{en conséq.} : si tout objet de $C$ est $\varprojlim$ d'objets de $A$, alors
le foncteur canonique
\[
r^{\circ} : C^{\circ} \longrightarrow \widehat{A^{\circ}} \quad (\text{resp. } \widehat{A^{\circ}}_{ab})
\]
vers les foncteurs $\underline{\mathrm{Hom}}(A, \mathrm{Ens})$ (resp.\ $\underline{\mathrm{Hom}}(A, \mathrm{Ab})$)
\emph{covariants} sur $A$ à valeurs dans $\mathrm{Ens}$ (resp.\ $\mathrm{Ab}$) est pl.\ fidèle. Si de plus tout
objet de $C$ est $\varprojlim$ de certain type d'objets de $C$, alors l'image essentielle est contenue dans la
sous-cat.\ pleine de $\widehat{A^{\circ}}$ (resp.\ $\widehat{A^{\circ}}_{ab}$) formée des foncteurs
$A \to \mathrm{Ens}$ (resp.\ $A \to \mathrm{Ab}$) qui sont $\varprojlim$ des deux types de foncteurs
représentables — et l'image essentielle est \emph{égale} à cette dernière catégorie si

\page{31}
$C$ est stable par les limites projectives du type envisagé.

En particulier, on trouve

\emph{Prop.} Soient $C$ une cat.\ (resp.\ cat.\ additive), $A$ une sous-cat.\ pleine, $\mathcal{T}$,
$\mathcal{T}'$ deux sous-\add{ens.} \struck{catégories pleines} de $\mathrm{Ob}(\mathrm{Cat})$, supposons que $C$
soit stable par $\varprojlim$ de type $\mathcal{T}$ et $\varinjlim$ de type $\mathcal{T}'$, et que tout objet de
$C$ soit une $\varprojlim$ \add{dans $C$} de type $\mathcal{T}$ d'objets de $A$, et une $\varinjlim$ \add{dans $C$}
de type $\mathcal{T}'$ d'objets de $A$. Alors $C$ est équivalente à la fois à la sous-cat.\ pleine de
$\widehat{A}$ ($= \underline{\mathrm{Hom}}(A^{\circ}, \mathrm{Ens})$) (resp.\ $\widehat{A}_{ab}$ $=
\underline{\mathrm{Hom}}(A^{\circ}, \mathrm{Ab})$) formée des foncteurs $A^{\circ} \to \mathrm{Ens}$ (resp.\
$A^{\circ} \to \mathrm{Ab}$) qui sont $\varprojlim$ \struck{de type $\mathcal{T}$} de type $\mathcal{T}$ de
foncteurs repr., et à la \add{cat.\ opposée de la} sous-catégorie pleine de $\widehat{A^{\circ}} \simeq
\underline{\mathrm{Hom}}(A, \mathrm{Ens})$ (resp.\ $\widehat{A^{\circ}}_{ab} \simeq
\underline{\mathrm{Hom}}(A, \mathrm{Ab})$) qui sont \struck{objets des} $\varinjlim$ de type $\mathcal{T}'$ de
foncteurs représentables.\note{sur les pages 28 à 33, les flèches sous les $\lim$ sont celles de la page, là où
elles se lisent ; sur cette page les deux sens figurent, $\mathcal{T}$ avec la flèche vers la gauche,
$\mathcal{T}'$ avec la flèche double}

\page{32}
En fait, on trouve un diagr.\ ess.\ commutatif

\begin{tikzcd}[column sep=small, row sep=normal]
 & A \subset C \arrow[dl, "\mathrm{id}"'] \arrow[d, "\rho\ \wr"] \arrow[dr, "\mathrm{id}"] \arrow[drr, "\rho'\ \wr"] & & \\
A \subset \widehat{A}_{\mathcal{T}} \subset \widehat{A} & A \arrow[r, hook] & (\widehat{A^{\circ}}_{\mathcal{T}'})^{\circ} \arrow[r, hook] & (\widehat{A^{\circ}})^{\circ}
\end{tikzcd}

\note{nous redessinons : sur la page, les quatre flèches partent de $A \subset C$ vers la ligne du bas ; celle
marquée $\rho\ \wr$ aboutit à $\widehat{A}_{\mathcal{T}}$, celle marquée $\rho'\ \wr$ à
$(\widehat{A^{\circ}}_{\mathcal{T}'})^{\circ}$}

(resp.\ avec $\widehat{A}_{ab}$ et $\widehat{A^{\circ}}_{ab}$). On trouve en particulier une équivalence
\uncertain{remarquable}
\[
\widehat{A}_{\mathcal{T}} \simeq (\widehat{A^{\circ}}_{\mathcal{T}'})^{\circ}
\]
qui \uncertain{exige} d'être définie \uncertain{intrinsèquement} \uncertain{uniquement} en termes de $A$,
$\mathcal{T}$, $\mathcal{T}'$, \uncertain{i.e.} chaque fois qu'on a une sous-cat.\ $D$ de $\widehat{A}$ (resp.\
une sous-cat.\ additive de $\widehat{A}_{ab}$), stable par les $\varprojlim$ de type $\mathcal{T}$, et \add{dans
$\widehat{A}$} \uncertain{où} \add{on construit} les $\varinjlim$ de type $\mathcal{T}'$ (pas nécess.\ les mêmes
que dans $\widehat{A}$ !) \uncertain{et telle} que tout objet soit à la fois $\varprojlim$ \add{de type
$\mathcal{T}$} \add{dans $C$} de foncteurs repr., et $\varinjlim$ \add{de type $\mathcal{T}'$} \add{dans $C$} de
\uncertain{tels} foncteurs.

Si $A$ est réduit à un seul objet $A$, et (dans le cas resp.\ ($C$ additive)) $\mathcal{T} = \mathcal{T}' =$
ens.\ catégories discrètes finies, on trouve une \uncertain{antiéquivalence} \uncertain{entre} \ill{}
cat.\ modules libres de t.f.\ à gauche

\page{33}
et droite sur $k = \mathrm{End}(A)$, et \uncertain{la dualité} antiéquivalence n'est autre que la dualité de
tels modules.

Revenons à la condition où tout objet de $C$ est une $\varinjlim$ de $C$ de type $\mathcal{T}'$ d'objets de $A$.
Alors \uncertain{si} \ill{} catégorie $C'$ où les limites existent, si
$\underline{\mathrm{Hom}}_{\mathcal{T}'}(C, C')$ est la sous-catégorie pleine des
$\underline{\mathrm{Hom}}(C', C')$\note{sic} formée des foncteurs qui commutent \uncertain{aux} $\varinjlim$, on
trouve que
\[
\underline{\mathrm{Hom}}_{\mathcal{T}'}(C, C') \longrightarrow \underline{\mathrm{Hom}}(A, C')
\]
est pl.\ fidèle \struck{\ill{} \ill{} une équivalence de cat.\ si \ill{} $C$ est
$\varinjlim$}.\marginal{à vérifier}\note{le passage biffé l'est de deux traits, le second mot sous une
rature épaisse ; « à vérifier » est écrit dans la marge gauche, séparé du texte par un trait vertical}
Supposons $C$ additive, et $\mathcal{T}$ \uncertain{l'ens.} des cat.\ discrètes finies — donc tout objet de $C$
est somme d'objets de $A$, alors \uncertain{on trouve}, si $C'$ \uncertain{additive},
\[
\underline{\mathrm{Hom}}_{add}(C, C') \simeq \underline{\mathrm{Hom}}_{\frac{1}{2}add}(A, C') .
\]
Si \uncertain{on applique} ça à l'inclusion
\[
C \xrightarrow{\text{incl}} \underline{\mathrm{Hom}}_{\frac{1}{2}add}(A^{\circ}, \mathrm{Ab})
\]
\note{avant $\underline{\mathrm{Hom}}$, un $\widehat{A}_{ab}$ biffé}
et désignons par $D_C$\note{une lettre surchargée après $D_C$} l'image essentielle, on trouve donc

\page{34}
\[
\underline{\mathrm{Hom}}_{add}(D, C') \xrightarrow{\ \approx\ } \underline{\mathrm{Hom}}_{\frac{1}{2}add}(A, C')
\]
d'où un foncteur inverse
\[
\underline{\mathrm{Hom}}_{\frac{1}{2}add}(A, C') \xrightarrow{\ \approx\ } \underline{\mathrm{Hom}}_{add}(D, C')
\]
d'où un accouplement
\[
\begin{array}{cl}
D & \times\ \underline{\mathrm{Hom}}_{\frac{1}{2}add}(A, C') \longrightarrow C' \\
\cap & \\
\underline{\mathrm{Hom}}_{\frac{1}{2}add}(A^{\circ}, \mathrm{Ab}) &
\end{array}
\]
qu'on note par \uncertain{ex.}
\[
(F, G) \longmapsto F \otimes_{C} G
\]
[dans le cas où $A$ est réduit à un \uncertain{objet}, donc avec sa str.\ ½ add.\ est défini par un anneau $k$,
$F$ est un \add{$k$-}module à droite libre de t.f., $G$ est un objet de $C'$ où $k$ \uncertain{opère} (i.e.\ un
$k$-$C'$-module à gauche) et $F \otimes G$ est le produit tensoriel « ordinaire ».]\note{le crochet n'est pas
refermé sur la page}

Si on considère \uncertain{alors} le plongement dual
\[
C \xrightarrow{\ \approx\ } D'^{\circ} \subset \underline{\mathrm{Hom}}_{\frac{1}{2}add}(A, \mathrm{Ab})^{\circ} ,
\]
on trouve
\[
\underline{\mathrm{Hom}}_{add}(D'^{\circ}, C') \xrightarrow{\ \approx\ } \underline{\mathrm{Hom}}_{\frac{1}{2}add}(A, C')
\]
d'où
\note{sous la ligne, une flèche courbe revient de droite à gauche, marquée $\approx$ : le foncteur inverse, comme
en tête de page}

\page{35}
i.e.
\[
\begin{array}{cl}
D'^{\circ} & \times\ \underline{\mathrm{Hom}}_{\frac{1}{2}add}(A, C') \longrightarrow C' \\
\cap & \\
\underline{\mathrm{Hom}}_{\frac{1}{2}add}(A, \mathrm{Ab})^{\circ} &
\end{array}
\]
noté
\[
\underline{\mathrm{Hom}}_{C}(F, G) .
\]
\note{une ligne de tirets, puis un trait horizontal, séparent la suite}

On suppose maintenant qu'on \uncertain{a} une catégorie $\mathcal{S}$, une sous-catégorie pl.\ $C$ \add{de
$\mathcal{S}_{ab}$ stable par sommes directes} (\uncertain{donc} $C$ est additive) et une sous-catégorie pleine
$A$ de $C$ telle que tout objet de $C$ soit isomorphe \add{dans $C$ donc dans $\mathcal{S}_{ab}$} à une somme
directe \add{finie} d'objets de $A$.\note{les compléments entre crochets obliques sont des ajouts en interligne}

On a alors \struck{\ill{}}
\[
\mathcal{S}^{\circ} \longrightarrow \struck{\ill{}}\ \underline{\mathrm{Hom}}_{add}(C, \mathrm{Ab})
\]
\struck{\ill{}} en associant à $X \in \mathrm{Ob}\,\mathcal{S}$ le foncteur
$M \longmapsto \underline{\mathrm{Hom}}_{\mathcal{S}}(X, M)$.\note{le « $\mathcal{S}^{\circ}$ » est récrit sur une
autre lettre ; devant, un signe \uncertain{« $\alpha$ »}}
Or ce foncteur est \uncertain{connu} à équiv.\ près quand on connaît son image dans
$\underline{\mathrm{Hom}}_{\frac{1}{2}add}(A, \mathrm{Ab})$, \uncertain{donc} \uncertain{c'est} \ill{}
$C$ est équivalente à celle des

\page{36}
\[
\mathcal{S}^{\circ} \longrightarrow \underline{\mathrm{Hom}}_{\frac{1}{2}add}(A, \mathrm{Ab})
\]
Nous ferons cette étude dans le cas où $\mathcal{S} = \widehat{\Delta} = \mathrm{Ss}$, et dans chacun des
\uncertain{2 cas} suivants \struck{\ill{} \ill{}} :

1°)\marginal{$C$ abélienne}\note{la note marginale est reliée par un trait au « 1° »} $C$ est formé des groupes
ab.\ ½ simpliciaux \add{ab.\ $\partial$-parfaits} \struck{\ill{}}, $A$ est la sous-catégorie formée des
$\Phi(i,n)$ ($n \in \mathbf{N}^{*}$ \struck{\ill{}}, $i \geq 1$) et \uncertain{$\Lambda T_{\bullet}$}
$(j \geq 0)$ ;\note{le dernier symbole est surchargé ; lecture douteuse}
\uncertain{on note} $X_{\bullet} \in \mathrm{Ob}\,\widehat{\Delta}$ \uncertain{fixé}. On note que la
connaissance des $\underline{\mathrm{Hom}}_{\widehat{\Delta}}(X_{\bullet}, M_{\bullet})$ comme
\uncertain{foncteur} \uncertain{en} $M_{\bullet}$ \uncertain{est réduite} \ill{} \uncertain{covariant en}
$M_{\bullet}$ \ill{} \uncertain{à} celle de $\underline{\mathrm{Hom}}_{\widehat{\Delta}}(X_{\bullet}, \varphi)$ pour
$\varphi$ dans $A$. On peut \uncertain{même} \ill{} $\varphi$ dans $A_0$ \uncertain{défini} \uncertain{ci-dessous}
\struck{\ill{}} $\varphi$ \ill{} \ill{} \ill{}, \uncertain{et} \ill{} pour $\varphi$ \uncertain{parcourant}
les $\Phi^{i}$.

2°)\marginal{($C_0$ est additive mais non abélienne.)} $C_0$ \struck{\ill{}} formée des groupes ab.\ ½
simpliciaux \struck{\ill{}} \uncertain{apériodiques} (\add{i.e.}\ \struck{\ill{}} $\in C$ et les
$\pi_i(\ldots)$ \uncertain{sont} sans torsion), $A_0 \subset C_0$ sous-cat.\ de ceux qui sont de la forme
\[
\Phi^{i} = \dot{\Lambda}\Phi_{\bullet} \ (i \geq 1), \qquad \Psi^{j} = \Lambda T_{\bullet} \ (j \geq 0).
\]
\note{les deux dernières formules sont très surchargées (un symbole biffé après $\dot{\Lambda}\Phi$) ; les
seconds membres sont des lectures douteuses}

\page{37}
Il y a lieu de \uncertain{travailler}
1°) \uncertain{au cas où} \ill{} \ill{} un anneau \uncertain{de base} $k$ \uncertain{de} Dedekind,
\add{(\uncertain{év.} \ill{})}
\uncertain{et} 2°) avec \uncertain{la généralité} \ill{} (\uncertain{donc} $C \subset \widehat{\Delta}_{k\text{-}\mathrm{mod}}$).
La \uncertain{situation} \uncertain{conduisant} \uncertain{à} \uncertain{retenir}, c'est que \uncertain{la}
\uncertain{connaissance} \uncertain{des} \uncertain{homs de} $X_{\bullet}$ dans les $k$-modules \uncertain{sh.}
\add{parfaits} \struck{\ill{}} $M$ (comme foncteur \uncertain{de} $M$), \uncertain{elle qui} \uncertain{connaît}
le complexe \add{de cochaînes « \uncertain{non dég.} »} $C^{*}(X_{\bullet}, k)$ \uncertain{et} \ill{}
\uncertain{dans} $k$. \uncertain{Notre} \uncertain{propos} \uncertain{principal} \uncertain{est} \ill{}
\uncertain{comment} \uncertain{reconstituer} la structure multiplicative \ill{} « \ill{} \ill{} » sur
\uncertain{ce} foncteur
\[
M_{\bullet} \longmapsto F_{X_\bullet}(M_\bullet) = F(M_\bullet) = \underline{\mathrm{Hom}}_{\widehat{\Delta}}(X_\bullet, M_\bullet)
\]
i.e.\ \uncertain{les homs}
\[
\begin{cases}
F(M_\bullet) \otimes_{k} F(N_\bullet) \longrightarrow F(M_\bullet \otimes N_\bullet) \\
\Gamma^{i} F(M_\bullet) \longrightarrow F(\Gamma^{i}(M_\bullet))
\end{cases}
\]
\note{dans la seconde ligne, un symbole biffé à l'intérieur de la première parenthèse ; le $\Gamma^{i}$ initial se
lit sur un trait surchargé}
en tant que \struck{\ill{}} structures \add{de cochaînes} sur le complexe $C^{*} = C^{*!}(X_\bullet, k)$.
\note{deux traits verticaux dans la marge gauche, en regard de la dernière phrase}
Il faudrait examiner \uncertain{si} \uncertain{pour} le \uncertain{cas} \add{\uncertain{$C^{*}$ connu}}
\struck{\ill{}} structures \uncertain{supplém.}

\page{38}
est déterminée par le complexe de De Rham-Sullivan de $X_\bullet$, et dans le cas $k$ quelconque (\uncertain{et
cas} 2° \uncertain{disons}) — le principe \uncertain{est} \uncertain{que} 1°) il est déterminé par les complexes
de De Rham à p.d.\ \struck{Enfin,} \struck{il y a} (il est clair que $C^{*}$ avec les structures suppl.\
envisagées \add{\uncertain{rectifiant}} \struck{\ill{}} les complexes de De Rham-Sullivan resp.\ le complexe de
De Rham à p.d.) Enfin, il faut voir dans quelle mesure $C^{*}$ avec ses structures multiplicatives à p.d.\
permet de reconstituer le type d'homotopie du $X_\bullet$ dont il provient. (Enfin, \uncertain{bien sûr}
\uncertain{faisceautiser} et relativiser\ldots\ !!!)

\note{le reste du feuillet est blanc, à part deux traits verticaux dans la marge droite}

\page{39}
Soit \add{$k$ anneau \uncertain{commutatif}} $\mathbb{K}$ une catégorie \add{$k$-}additive (i.e.\ additive, munie
de $k \to \mathrm{End}(\mathrm{id}_{\mathbb{K}})$), $\mathbb{C}^{\bullet}$ un complexe \add{de cochaînes} dans
$\mathbb{K}$
\[
0 \to \mathbb{C}^{0} \xrightarrow{d^{0}} \mathbb{C}^{1} \xrightarrow{d^{1}} \cdots \to \mathbb{C}^{i}
\xrightarrow{d^{i}} \mathbb{C}^{i+1} \to \cdots
\]
On suppose :

a)
\[
\operatorname{Hom}(\mathbb{C}^{i}, \mathbb{C}^{j}) \simeq
\begin{cases}
0 & \text{si } j \neq i, i+1 \\
k & \text{si } j = i, \text{ base } \mathrm{id}_{\mathbb{C}^{i}} \\
k & \text{si } j = i+1, \text{ base } d^{i}
\end{cases}
\]

b) Tout objet de $\mathbb{K}$ est isomorphe à une somme \add{finie} d'objets de la forme $\mathbb{C}^{i}$.

Soit $\mathcal{A}$ la sous-catégorie pleine de $\mathbb{K}$ formée des \struck{\ill{}} $\mathbb{C}^{i}$ ; elle
n'est pas additive, mais hérite d'une structure $k$-linéaire, i.e.\ \struck{\ill{}} les
$\operatorname{Hom}_{\mathcal{A}}$ \uncertain{sont} \uncertain{des} $k$-modules, \uncertain{la} composition
\uncertain{est} $k$-bil.

\struck{E} L'existence de $(\mathbb{K}, \mathbb{C}^{\bullet})$ satisfaisant a) et b) est immédiate par
construction « abstraite ». On peut aussi réaliser $\mathbb{K}$ comme \add{la} cat.\ \add{des} complexes de
chaînes \add{bornés} de $k$-modules satisfaisant les conditions
\begin{enumerate}
\item[1°)] \struck{\ill{}} $L_\bullet$ est acyclique i.e.\ $H_i(L_\bullet) = 0$ $\forall\, i \geq 0$
\item[2°)] $L_n$ et $Z_n(L_\bullet)$ sont libres de t.f.\ pour tout $n$
\item[3°)] $L_n$ est nul pour $n$ grand.
\end{enumerate}
On prendra alors
\[
\mathbb{C}^{i} = C^{i}_{\bullet} = \lbrace \cdots 0 \to 0 \to
\underset{\text{degré } i+1}{k} \xrightarrow{\ \mathrm{id}\ } \underset{\text{degré } i}{k} \to 0 \to \cdots \to 0 \rbrace
\]
\note{au-dessus du « 0 » de droite, « degré 0 » ; avant $C^{i}_{\bullet}$, un premier membre $\mathbb{K}(\ldots)$ est biffé}
i.e.\ les comp.\ sont nulles sauf en degré $i$, $i+1$ \uncertain{où} elles sont $k$, et l'opér.\ diff.\
\add{$d_{i+1}$} \struck{\ill{}} est l'identité. On prend $d^{i} : \mathbb{C}^{i} \to \mathbb{C}^{i+1}$
\struck{\ill{}} caractérisé par $(d^{i})_{i+1} = \mathrm{id}_k$. Les conditions a) et b) sont satisfaites
trivialement. On pourrait aussi \uncertain{travailler} \add{\uncertain{dans}} des complexes de cochaînes : le
plus,
\note{dans toute la marge gauche, une colonne de petits complexes dessinés verticalement (flèches $d$, $0$,
accolades, un $\Sigma$), avec les mentions obliques « degré $i+1$ », « degré $i$ », « degré 0 », qui illustrent
les $C^{i}_{\bullet}$ ; une ligne de ce croquis rejoint le bas de la page. Non redessiné}

\page{40}
\note{la moitié supérieure de la page est barrée de trois longues diagonales ; elle se lit et est transcrite en
entier}
[$\mathcal{C}$ catégorie] des $k$-modules ½ simpliciaux \add{$L_*$} \struck{\ill{}} qui sont flasques et str.\
$\partial$-parfaits, i.e.\ qui satisfont \struck{aux} \struck{deux} conditions équivalentes
\begin{enumerate}
\item[a)] $L_*$ est \struck{\ill{}} \add{isomorphe} à une somme directe finie de modules ½ simpliciaux
$C^{i}_{*} = \Lambda^{i+1}\Phi_{*k}$ $(i \geq 0)$\note{sous cette formule, une seconde biffée et surchargée,
\uncertain{$\Phi(i)_{*} = \Lambda^{i}\Phi_{*} \ldots$ $(i \geq 1)$}, avec « $= C^{i-1}_{*}$ » au-dessous}
\item[b)] Les $L_n$ sont libres de type fini, de rang constant sur $\operatorname{Spec} k$, la fonction
$n \mapsto \operatorname{rg}_k L_n$ est polynomiale pour $n$ grand, $L_*$ est flasque i.e.\
$\pi_i(L_*) = 0$ pour $i \geq 0$.
\end{enumerate}
\marginal{On suppose $k$ \uncertain{anneau} \uncertain{de} \uncertain{Dedekind} (\uncertain{proj.}\ de type
fini sur $k$ \uncertain{libres} par le principal)}\note{note oblique dans la marge gauche, en regard de a) ;
lecture très incertaine}

$\mathcal{A} \subset \mathcal{C}$ la sous-catégorie pleine formée des \struck{\ill{}} $C^{i}_{*}$.

\emph{NB} $\mathcal{C}$ est additive (et $k$-\uncertain{linéaire}), $\mathcal{A}$ n'est pas stable par sommes
finies, mais hérite une structure $k$-linéaire. On a
\[
\underline{\mathrm{Hom}}(C^{i}_{*}, C^{j}_{*}) =
\begin{cases}
0 & \text{si } j \neq i, i+1 \\
k & \text{si } j = i, \text{ base } \mathrm{id}_{C^{i}_{*}} \\
k & \text{si } j = i+1, \text{ base } d^{i} = T_{*}\wedge
\end{cases}
\]
Notons que \struck{\ill{} $C^{i}_{*} \ldots C^{i}_{*}$ \ill{}}\note{une ligne entière biffée et
surchargée}
$C^{\bullet}_{*} = (C^{i}_{*}, T_{*}\wedge)$ est un complexe de cochaînes dans $\mathcal{A}$ ou dans
$\mathcal{C}$ ou dans $\widehat{\Delta}_{k\text{-lin}}$ — il jouera le rôle du complexe de cochaînes
\emph{universel}.

\struck{\ill{}}

\emph{Prop 1} Soit $\mathcal{B}$ une catégorie \add{$k$-additive} \struck{additive}. Alors
\[
\underline{\mathrm{Hom}}(\mathcal{C}, \mathcal{B}) \xrightarrow[\ \approx\ ]{\text{rest}}
\underline{\mathrm{Hom}}_{k\text{-lin}}(\mathcal{A}, \mathcal{B})
\]
\note{sous le premier $\underline{\mathrm{Hom}}$, l'indice « $k$-lin » est biffé}
et, par $F \mapsto F(C^{\bullet}_{*})$, le second est $\approx$ à
$\mathcal{B}^{\bullet}_{k}$ = cat.\ des \add{complexes de} cochaînes \add{de $\mathcal{B}$}.\note{« complexes
de » est entouré}

\struck{\ill{} Soit $K^{\bullet} \in \mathrm{Ob}\,\mathcal{B}^{*}$ un complexe de cochaînes de $\mathcal{B}$}

\emph{Exemple 1} [\add{\uncertain{Faisons} $\mathcal{B} = \mathit{Ab}_k$.}] Soit $L_\bullet \in \mathrm{Ob}\,\mathcal{C}$,
considérons le foncteur cov.\ qu'il représente
\[
\mathcal{C} \longrightarrow \mathit{Ab}_k, \qquad M_* \longmapsto \operatorname{Hom}_{\mathcal{C}}(L_*, M_*) ;
\]
on a donc un foncteur pl.\ fid.\ (\uncertain{vu} \add{\uncertain{sur} des foncteurs repr.}\ \ldots)
\[
\mathcal{C}^{\circ} \hookrightarrow \underline{\mathrm{Hom}}_{k\text{-lin}}(\mathcal{C}, \mathit{Ab}_k)
\xrightarrow{\ \approx\ } \underline{\mathrm{Hom}}_{\mathrm{ac}}(\mathcal{A}, \mathit{Ab}_k)
\xrightarrow{\ \approx\ } \mathit{Ab}_k^{\bullet}
\]
\note{l'indice du dernier $\underline{\mathrm{Hom}}$ se lit « ac » ou « $k$-lin », lecture douteuse}
\note{la page 41 (lot 3) reprend : « On trouve donc un foncteur pl.\ fid.\ $\mathcal{C}^{\circ}
\hookrightarrow \mathit{Ab}_k^{\bullet}$ »}

\end{document}
